%% ============================================================
%%  A3 -- Phase 4 Signal Research: BTCUSDT Perpetual Futures
%%  Compile with: pdflatex A3.tex  (twice for refs)
%% ============================================================
\documentclass[12pt, a4paper]{article}

\usepackage{amsmath, amssymb, amsthm, mathtools}
\usepackage{booktabs}
\usepackage{hyperref}
\usepackage{geometry}
\usepackage{xcolor}
\usepackage{array}
\usepackage{multirow}
\usepackage{enumitem}
\usepackage{caption}
\usepackage{float}
\usepackage{parskip}
\usepackage{bm}

\geometry{margin=2.5cm}

\hypersetup{
  colorlinks=true,
  linkcolor=blue!60!black,
  urlcolor=blue!60!black
}

\theoremstyle{plain}
\newtheorem{theorem}{Theorem}[section]
\newtheorem{proposition}{Proposition}[section]
\newtheorem{corollary}{Corollary}[section]
\theoremstyle{definition}
\newtheorem{definition}{Definition}[section]
\newtheorem{remark}{Remark}[section]

\newcommand{\IC}{\mathrm{IC}}
\newcommand{\E}{\mathbb{E}}
\newcommand{\Var}{\mathrm{Var}}
\newcommand{\Cov}{\mathrm{Cov}}
\newcommand{\Cor}{\mathrm{Cor}}
\newcommand{\DVOL}{\mathrm{DVOL}}
\newcommand{\bp}{\,\mathrm{bp}}
\newcommand{\OI}{\mathrm{OI}}

\title{\textbf{Phase 4 Signal Research --- BTCUSDT Perpetual Futures}\\[0.4em]
\large Mark-Index Basis, Cross-Asset Implied Volatility, and Liquidation Exhaustion\\[0.3em]
\large A Complete Treatment from First Principles}
\author{}
\date{May 2026}

\begin{document}
\maketitle
\tableofcontents
\newpage

%% ============================================================
\section{Introduction}
%% ============================================================

This document records Phase 4 of the BTCUSDT perpetual futures signal
research programme.  It is a self-contained technical record covering
the design, derivation, and empirical validation of three genuinely new
signal families, each based on a distinct economic mechanism.

The core result is: \textbf{Q5c (BTC-excess implied volatility) advances
to full validation with OOS Sharpe $+3.30$ ($n = 105$ trades,
block-bootstrap $p < 0.0005$), is independent of N3 (incremental IC
ratio $2.70\times$, $\Cor = +0.342$), and subsequently superseded when
the research programme pivoted to higher-frequency signal exploration
in Phase 5.}  All other Phase 4 signals are killed or untestable.

The validated signal stack entering Phase 4 is:
\begin{enumerate}
  \item N3 (DVOL $z$-score $\geq 0.75$, DVOL $\geq 54$): OOS Sharpe $+2.95$,
        $1.2$ trades/week (\textbf{A1}).
  \item P3 (OI-Price Divergence, DD regime, DVOL $\geq 54$): OOS Sharpe $+5.18$
        on exclusive trades, independent of N3 (\textbf{A2}).
\end{enumerate}

Phase 4 tests three genuinely new mechanisms, none overlapping with
previously killed signals:

\begin{description}
  \item[Q4 --- Mark-Index Basis]
        The instantaneous premium or discount of the perpetual mark price
        relative to the Binance spot index, measured continuously rather
        than averaged over 8-hour funding windows.

  \item[Q5 --- ETH-BTC Cross-DVOL]
        Three variants exploiting the cross-asset relationship between
        Deribit BTC and ETH implied volatility: ETH $z$-score alone (Q5a),
        the product of both $z$-scores (Q5b), and the difference
        BTC$_{N3z}$ $-$ ETH$_{N3z}$ (Q5c, the ``BTC-excess'' signal).

  \item[Q6 --- Liquidation Exhaustion]
        A spike in daily forced long-liquidation notional marks the
        final stage of a leveraged sell-off.  Data source: Binance Data
        Vision daily liquidation snapshots.
\end{description}

The report proceeds as follows. Section~\ref{sec:market} introduces
the relevant market structures from first principles: the mark-index
basis and funding mechanics, cross-asset implied volatility and spillover
theory, and crypto liquidation cascade dynamics.
Section~\ref{sec:stats} summarises the statistical framework.
Sections~\ref{sec:q4}--\ref{sec:q6} analyse each signal.
Section~\ref{sec:discussion} provides a synthesis.
Section~\ref{sec:conclusion} concludes.

%% ============================================================
\section{Market Background}
\label{sec:market}
%% ============================================================

\subsection{The Perpetual Mark Price and the Basis}
\label{sec:basis_theory}

\subsubsection{Mark Price Construction}

Binance computes the perpetual mark price $M_t$ as a function of the
spot index $I_t$ and a smoothed premium:
\begin{equation}
  M_t = I_t \times \left(1 + \mathrm{FairBasis}_t\right)
  \label{eq:mark}
\end{equation}
where $\mathrm{FairBasis}_t$ is an exponentially weighted average of
recent premium index observations.  The \emph{premium index} at each
minute is:
\begin{equation}
  P_t^{\text{idx}} = \frac{\max(B_t, I_t) - \min(A_t, I_t)}{I_t}
  - \max(A_t, I_t) \cdot \frac{r_{\text{int}}}{8} / I_t
  \label{eq:premidx}
\end{equation}
where $B_t$ is the best bid price of the perpetual, $A_t$ is the best
ask price, and $r_{\text{int}} = 0.03\%$ is the 8-hour interest rate.
The second term is a small interest adjustment.  In practice,
$P_t^{\text{idx}} \approx (\bar{P}_t - I_t) / I_t$, where $\bar{P}_t$
is the mid-price of the perpetual.

\subsubsection{The Basis and its Relationship to Funding}

\begin{definition}[Mark-Index Basis]
  The \emph{basis} at time $t$ is the relative premium of the perpetual
  mark price over the spot index:
  \begin{equation}
    b_t = \frac{M_t - I_t}{I_t}
    \label{eq:basis}
  \end{equation}
  A positive basis means the perpetual trades above spot (net long
  demand); a negative basis means the perpetual trades below spot
  (net short demand or forced long selling).
\end{definition}

The basis and the 8-hour funding rate are directly linked.  The funding
rate at settlement time $\tau$ (every 8 hours) is approximately:
\begin{equation}
  F_\tau \approx \frac{1}{480}\sum_{t=\tau-480}^{\tau} P_t^{\text{idx}}
               \approx \bar{b}_{[\tau-480,\tau]}
\end{equation}
where the average is over the 480 minutes in the preceding settlement
window.  This means:

\begin{itemize}
  \item A persistently positive basis in the settlement window implies
        a positive funding rate: longs will pay shorts.
  \item A persistently negative basis implies a negative funding rate:
        shorts will pay longs.
\end{itemize}

The key difference from the H1 funding signal (killed in Phase 1) is
resolution: the 8-hour funding rate is a \emph{lagging average} of the
basis over the preceding window.  The basis $b_t$ is the \emph{real-time}
instantaneous measure.

\subsubsection{Two Competing Hypotheses}

The basis generates two economically opposite predictions:

\begin{description}
  \item[Q4a Momentum:] A positive basis (perpetual at premium to spot)
        signals excess long demand.  Longs are willing to hold and pay the
        upcoming funding, indicating conviction.  The premium persists and
        predicts positive 24h returns.

  \item[Q4b Contrarian:] A negative basis (perpetual at discount to spot)
        signals forced selling by longs who can no longer bear the cost of
        a negative funding rate.  Once this forced selling clears, the
        discount reverses and predicts positive 24h returns.
\end{description}

These hypotheses make opposite sign predictions for the IC.  Empirically
testing both is standard practice: the data determines which mechanism
dominates.

\subsubsection{Descriptive Statistics}

The basis series is computed from the Binance FAPI premium index klines
(\texttt{indexPriceLine} and \texttt{markPrice} 1-minute bars) over the
2023-01-01 to 2026-05-13 OOS window:

\begin{center}
\begin{tabular}{lrrrr}
\toprule
Statistic & Mean (bp) & Std (bp) & Min (bp) & Max (bp) \\
\midrule
OOS daily basis & $-2.62$ & $3.51$ & $-7.12$ & $+12.21$ \\
\bottomrule
\end{tabular}
\end{center}

The mean daily basis of $-2.62\bp$ reflects the persistent longs-paying
environment (positive average funding rate) during the sample: when
longs pay funding, the equilibrium basis drifts slightly negative as
perpetual prices are dragged toward spot.  The range of $[-7.12,
+12.21]\bp$ encompasses both deep discounts (acute forced long selling)
and large premiums (euphoric long demand).

\subsection{Cross-Asset Implied Volatility and Spillover}
\label{sec:crossvol_theory}

\subsubsection{The Deribit ETH DVOL}

Just as Deribit publishes a BTC 30-day implied volatility index
(BTCVOL / DVOL, described fully in A1 Section 2.2), it publishes an
equivalent index for ETH options:

\begin{definition}[ETH DVOL]
  The Deribit ETH implied volatility index (\texttt{ETHDVOL}) is
  constructed by the same model-free VIX methodology (Britten-Jones
  \& Neuberger, 2000) applied to ETH/USD options on Deribit, targeting
  a 30-calendar-day constant maturity.  It is published at hourly
  frequency in annualised percentage units:
  \[
    \DVOL^{\mathrm{ETH}}_t = 100 \times \sqrt{\hat{\sigma}^2_{30\mathrm{d}}(\mathrm{ETH}) \times 365}
  \]
\end{definition}

The two indices share a common fear factor (broad crypto market
uncertainty) but capture distinct liquidation dynamics: ETH has its own
leverage ecosystem, DeFi protocol liquidations, and smart-contract-specific
risks that can cause ETH vol to spike independently of BTC vol.

Over the sample period (2022--2026, 38{,}281 hourly observations):
\begin{itemize}[noitemsep]
  \item Mean BTC DVOL: 52.4\%; mean ETH DVOL: 58.1\%.
  \item Pearson correlation of the two raw series: $+0.82$.
  \item Pearson correlation of the two 30-day $z$-scores: $+0.696$.
\end{itemize}

The high raw correlation ($+0.82$) reflects the common macro fear factor.
The lower $z$-score correlation ($+0.696$) reflects that the two assets
have distinct short-term deviation patterns around their own 30-day
histories --- the component exploited by Q5c.

\subsubsection{Volatility Spillover Theory}

Diebold \& Yilmaz (2012) quantify directional volatility spillovers
across financial markets using a forecast error variance decomposition
of a vector autoregression.  Their finding is that implied volatility
in one asset predicts future realised volatility in correlated assets
--- a ``spillover'' from the fear of one market to the price dynamics of
another.

In crypto, BTC and ETH are the two dominant assets.  They share a
common speculative demand factor (``crypto beta''), meaning:
\begin{itemize}[noitemsep]
  \item A broad crypto fear event tends to spike both BTC and ETH DVOL
        simultaneously.
  \item A BTC-specific event (e.g., a large miner capitulation, a
        Bitcoin ETF approval) may spike BTC DVOL without equally affecting
        ETH DVOL.
  \item An ETH-specific event (e.g., a DeFi hack, a major protocol
        upgrade) may spike ETH DVOL without equally affecting BTC DVOL.
\end{itemize}

This distinction motivates three signal variants:

\begin{description}
  \item[Q5a:] ETH N3$z$ alone predicts BTC returns (ETH fear leads BTC
        price recovery).
  \item[Q5b:] The product BTC$_{N3z} \times$ ETH$_{N3z}$ captures joint
        confirmation (both assets simultaneously fearful signals a broad
        crypto stress event with a strong mean-reversion).
  \item[Q5c:] The difference BTC$_{N3z} -$ ETH$_{N3z}$ captures
        BTC-specific fear elevation above the crypto baseline.  When BTC's
        $z$-score substantially exceeds ETH's, BTC-specific forced selling
        has exceeded systemic fear.  The hypothesis is that this excess
        reverts in the 24-hour window.
\end{description}

\subsubsection{The BTC-Excess Mechanism (Q5c)}

The Q5c signal deserves the most detailed treatment as it is the Phase 4
advance.  The economic story:

\begin{enumerate}
  \item A BTC-specific event causes BTC DVOL to spike.  For example:
        a large BTC holder capitulates, triggering a cascade of long
        liquidations specific to BTC perpetuals.
  \item ETH DVOL rises less because the event is BTC-specific.
        $N3z^{\mathrm{BTC}} \gg N3z^{\mathrm{ETH}}$.  Q5c is large and positive.
  \item The BTC-specific forced supply exhausts (the large holder
        has sold, the cascade has cleared, over-levered longs are gone).
  \item BTC price recovers from the oversold level over the next 24 hours.
        ETH recovers less (it did not fall as much).
  \item Q5c predicts positive 24h BTC returns with IC $+0.189$.
\end{enumerate}

The mechanism is a refinement of the fear-resolution channel (A1 Section
2.3) and the liquidation-exhaustion channel (A2 Section 2.1): by
conditioning on the BTC-ETH spread, Q5c filters out broad crypto sell-offs
(where both assets fall similarly and the recovery is less predictable)
in favour of BTC-specific overshoots.

\subsection{Liquidation Cascade Dynamics}
\label{sec:cascade_theory}

\subsubsection{How Cascades Form}

A \emph{liquidation cascade} occurs when forced selling by one set of
over-levered participants causes a price decline large enough to trigger
additional forced selling by a second set:

\begin{enumerate}
  \item Participant A holds a $10\times$ levered BTC long.  The mark price
        falls 9\%, their margin is below maintenance.  The exchange's
        liquidation engine sells their position at market.
  \item The forced sale pushes the mark price lower.  Participant B,
        also levered long, now breaches their maintenance margin.
        Their position is also liquidated.
  \item The cascade continues until all positions above a certain
        leverage level are wiped out.  OI drops sharply.  The mark
        price has fallen further than the underlying ``fair value'' would
        justify, creating a temporary oversold condition.
\end{enumerate}

Daily forced long-liquidation notional is the most direct observable of
cascade severity.  A spike ($>2\sigma$ above the 30-day mean) indicates
a major cascade has occurred, and the hypothesis is that the subsequent
24h return is strongly positive as the oversold condition resolves.

\subsubsection{Data Availability}

The Binance Data Vision platform publishes daily liquidation snapshot
files at:
\begin{center}
\texttt{/data/futures/um/daily/liquidationSnapshot/BTCUSDT/}
\end{center}
These files contain all individual liquidation orders for a given day,
from which daily long-liquidation notional can be aggregated.

During Phase 4 data collection, this endpoint returned HTTP 404
for all dates 2022--2026, rendering Q6 untestable from this source.
Alternative sources (Coinalyze, CoinGlass, Velo Data) provide similar
data but require paid subscriptions or have incomplete historical
coverage; these were not available for this research phase.  Q6 is
formally recorded as \textbf{untestable} in Phase 4.

%% ============================================================
\section{Statistical Framework}
\label{sec:stats}
%% ============================================================

This section derives the statistical tools used throughout Phase 4 from
first principles.  The treatment is self-contained.

\subsection{Log-Returns as Prediction Targets}

For a price series $\{C_t\}$, the \emph{simple return} over a horizon $h$ is
$\tilde{r}_t^{(h)} = C_{t+h}/C_t - 1$.  The \emph{log-return} is:
\[
  r_t^{(h)} = \log\frac{C_{t+h}}{C_t} = \log C_{t+h} - \log C_t
\]
Log-returns are preferred for three reasons:

\paragraph{Temporal additivity.}
Log-returns over sub-periods add exactly: $r_{t,t+h+k} = r_{t,t+h} + r_{t+h,t+h+k}$.
Simple returns compound multiplicatively, requiring chain multiplication.
This additivity is essential for equity curve construction: the log-equity
curve is the cumulative sum of log-returns.

\paragraph{Symmetry.}
A simple return of $+50\%$ is not cancelled by a subsequent $-50\%$:
$1.5 \times 0.5 = 0.75 \neq 1$.  Log-returns are symmetric: $+\log 1.5$
is exactly cancelled by $-\log 1.5$.  This symmetry makes long and short
positions analytically equivalent, since a short log-return is $-r_t^{(h)}$.

\paragraph{Cost comparability.}
Transaction costs are approximately constant fractions of notional regardless
of direction.  A round-trip cost $c_{\text{RT}}$ is subtracted uniformly from
$r_t^{(h)}$ regardless of its sign or magnitude.  This is exact in log-return
space and approximate in simple return space.

\paragraph{Basis-point convention.}
All PnL, drawdown, and per-trade figures are stated in \emph{basis points}:
$1\bp = 0.01\% = 10^{-4}$.  A log-return of $0.01$ (1\%) equals $100\bp$.

The 24-hour total return including funding payments is:
\begin{equation}
  \pi_t^{(24h)} = \log\frac{C_{t+1440}}{C_t}
                  - \sum_{k:\,\tau_k \in [t, t+1440]} F_{\tau_k}
\end{equation}
where $F_{\tau_k}$ is the 8-hour funding rate at settlement time $\tau_k$,
and the sum covers the three settlements within the hold window.
Round-trip maker cost: $c_{\mathrm{RT}} = 6\bp$ ($3\bp$ entry, $3\bp$ exit).

\subsection{The Spearman Information Coefficient}
\label{sec:IC}

\begin{definition}[Spearman IC]
  Let $\{(s_t, r_t)\}_{t=1}^n$ be a sample of signal-return pairs.
  Define the rank functions:
  \[
    R_t = \rank(s_t) = \left|\{j : s_j \leq s_t\}\right|, \qquad
    Q_t = \rank(r_t) = \left|\{j : r_j \leq r_t\}\right|
  \]
  (ranks from 1 to $n$; midranks for ties).  The Spearman IC is the Pearson
  correlation applied to the rank vectors:
  \begin{equation}
    \IC(s, r) = \frac{\sum_{t=1}^n (R_t - \bar{R})(Q_t - \bar{Q})}
                     {\sqrt{\sum_{t=1}^n (R_t - \bar{R})^2}
                      \sqrt{\sum_{t=1}^n (Q_t - \bar{Q})^2}}
    \label{eq:spearman}
  \end{equation}
  For continuous distributions with no ties, $\bar{R} = \bar{Q} = (n+1)/2$
  and $\sum (R_t - \bar{R})^2 = n(n^2-1)/12$, giving the familiar formula:
  \[
    \IC(s, r) = 1 - \frac{6\sum_{t=1}^n d_t^2}{n(n^2 - 1)},
    \qquad d_t = R_t - Q_t
  \]
\end{definition}

\paragraph{Why Spearman over Pearson?}
The Pearson correlation $\Cor(s, r)$ is the natural measure of
\emph{linear} predictability.  However:
\begin{enumerate}[noitemsep]
  \item Log-return distributions are leptokurtic; a few large-return days
        dominate the Pearson correlation.  A single outlier can flip the sign.
  \item The relationship between a signal and $r_t^{(24h)}$ need not be
        linear.  The Spearman IC is invariant to any monotone transformation
        of either variable.
  \item Trading strategies based on signal \emph{rank} (enter when the signal
        is in the top quintile) are naturally assessed by rank correlation,
        not level correlation.
\end{enumerate}

\begin{remark}
  A negative IC is equally informative as a positive IC of the same
  magnitude: the signal should simply be traded on the opposite side.
  We report $|\IC|$ and the sign separately throughout.
\end{remark}

\subsection{The Breakeven Information Coefficient}
\label{sec:ICstar}

A statistically significant IC does not imply profitability.  The signal must
generate sufficient expected gross profit to cover round-trip transaction costs.
We now derive the minimum IC required.

\begin{proposition}[Breakeven IC]
  \label{prop:ICstar}
  Consider a trading strategy that enters long when a standardised signal
  $z_t \sim \mathcal{N}(0,1)$ is positive and short when $z_t$ is negative,
  exiting after $h$ bars.  Under the linear-IC model
  \begin{equation}
    r_t^{(h)} = \IC \cdot \sigma_r \cdot z_t + \varepsilon_t, \qquad
    \varepsilon_t \perp z_t,\quad
    \Var(r_t^{(h)}) = \sigma_r^2
    \label{eq:linearIC}
  \end{equation}
  the expected gross profit per trade equals
  $\IC \cdot \sigma_r \cdot \sqrt{2/\pi}$.  The signal is profitable
  against a round-trip cost $c_{\mathrm{RT}}$ if and only if
  \begin{equation}
    |\IC| > \IC^* \;\coloneqq\; \frac{c_{\mathrm{RT}}}{\sigma_r \sqrt{2/\pi}}
    \label{eq:ICstar}
  \end{equation}
\end{proposition}

\begin{proof}
  Under model~\eqref{eq:linearIC}, since $z_t \perp \varepsilon_t$:
  \[
    \E[r_t^{(h)} \mid z_t] = \IC \cdot \sigma_r \cdot z_t
  \]
  For a long trade conditional on $z_t > 0$:
  \begin{align*}
    \E[r_t^{(h)} \mid z_t > 0]
    &= \IC \cdot \sigma_r \cdot \E[z_t \mid z_t > 0]
    = \IC \cdot \sigma_r \cdot \frac{\E[z_t \cdot \mathbf{1}_{z_t > 0}]}
                                     {P(z_t > 0)}
  \end{align*}
  For $z_t \sim \mathcal{N}(0,1)$:
  \[
    \E[z_t \mid z_t > 0]
    = \frac{\int_0^\infty z\,\phi(z)\,\mathrm{d}z}{\tfrac{1}{2}}
    = 2\int_0^\infty z\,\phi(z)\,\mathrm{d}z
    = 2\left[-\phi(z)\right]_0^\infty
    = 2\phi(0)
    = \sqrt{\tfrac{2}{\pi}}
  \]
  By symmetry, a short trade ($z_t < 0$) earns $\E[-r_t^{(h)} \mid z_t < 0]
  = \IC \cdot \sigma_r \cdot \sqrt{2/\pi}$ as well.

  Expected net profit per trade equals gross minus round-trip cost:
  \[
    \E[\text{net profit}] = \IC \cdot \sigma_r \sqrt{2/\pi} - c_{\mathrm{RT}}
  \]
  Setting this to zero and solving for IC gives~\eqref{eq:ICstar}.\qed
\end{proof}

\begin{remark}[Threshold entry]
  If the strategy only trades when $|z_t| > c$ for some $c > 0$, the expected
  gross profit per trade is $\IC \cdot \sigma_r \cdot \E[z \mid z > c]$.  Since
  $\E[z \mid z > c] = \phi(c)/[1-\Phi(c)] > \sqrt{2/\pi}$ for all $c > 0$,
  the breakeven IC is \emph{lower} than~\eqref{eq:ICstar}.  The formula as
  stated is conservative: it applies to the strategy that enters on every bar.
\end{remark}

\begin{remark}[Pearson vs.\ Spearman]
  Proposition~\ref{prop:ICstar} is stated for the Pearson IC.  The Spearman
  IC satisfies $\IC_S \approx (6/\pi)\arcsin(\IC_P / 2)$, so
  $\IC_S \approx \IC_P$ for $|\IC| \ll 1$.  Since all empirical ICs in this
  report are in $[-0.20, +0.20]$, the same $\IC^*$ formula is applied.
\end{remark}

For the parameters used throughout Phase 4:
\begin{itemize}[noitemsep]
  \item $c_{\mathrm{RT}} = 6\bp$ (maker: $3\bp$ entry, $3\bp$ exit).
  \item $\sigma_r^{(24h)} \approx 2.135\%$ (empirical OOS volatility).
  \item $\IC^*_{\mathrm{maker}} \approx 0.035$.
\end{itemize}

The IC/IC* ratio gives the ``multiple of breakeven'': a ratio of $3\times$
means the signal generates three times the expected gross profit needed
to cover costs.

\subsection{The Overlapping-Return Inflation Problem}
\label{sec:overlap}

The 1-minute BTCUSDT dataset has $T \approx 527{,}040$ observations in the OOS
2024 window.  If we compute $r_t^{(24\text{h})} = \log(C_{t+1440}/C_t)$
for every $t$ and treat all pairs as independent, we overstate statistical
significance severely.

\paragraph{The autocorrelation structure of overlapping returns.}
Under the assumption that 1-minute log-returns are i.i.d.\ with
variance $\delta^2$, the 24h log-return decomposes as:
\[
  r_t^{(24\text{h})} = \sum_{j=0}^{1439} \Delta_j(t),
  \qquad \Delta_j(t) = \log\frac{C_{t+j+1}}{C_{t+j}}
\]
Two adjacent observations share all but one minute of their window.
More generally, for lag $k \leq 1440$:
\begin{align*}
  \Cov\!\left(r_t^{(24\text{h})}, r_{t+k}^{(24\text{h})}\right)
  &= \Cov\!\left(\sum_{j=0}^{1439}\Delta_j(t),\,
                 \sum_{j=0}^{1439}\Delta_j(t+k)\right)
   = (1440 - k)\,\delta^2
\end{align*}
since the two sums share exactly $1440 - k$ common terms (for $k < 1440$).
The autocorrelation at lag $k$ is:
\begin{equation}
  \Cor\!\left(r_t^{(24\text{h})}, r_{t+k}^{(24\text{h})}\right)
  = \frac{(1440 - k)\,\delta^2}{1440\,\delta^2}
  = 1 - \frac{k}{1440}, \qquad 0 \leq k \leq 1440
  \label{eq:overlap_autocor}
\end{equation}
This is a tent function that linearly decays from 1 at lag 0 to 0 at lag 1440.

\paragraph{Effective sample size.}
For a stationary series with autocorrelation function $\rho_k$, the effective
sample size for variance estimation is approximately:
\[
  n_{\mathrm{eff}} \approx \frac{T}{1 + 2\sum_{k=1}^{\infty} \rho_k}
\]
Substituting~\eqref{eq:overlap_autocor}:
\[
  1 + 2\sum_{k=1}^{1440-1}\!\!\left(1 - \frac{k}{1440}\right)
  = 1 + 2 \cdot \frac{1440 \cdot 1439}{2 \cdot 1440}
  = 1 + 1439
  = 1440
\]
Therefore:
\begin{equation}
  n_{\mathrm{eff}} \approx \frac{T}{1440} \approx 366
  \label{eq:neff}
\end{equation}
The effective number of independent 24-hour return observations is
approximately 366 per year --- one per day --- regardless of whether we
compute the return at every 1-minute bar or sample once per day.

\paragraph{Practical implication.}
A $t$-statistic computed from the full $n = 527{,}040$ series is inflated
by a factor of $\sqrt{1440} \approx 38$ relative to the correct statistic
based on $n_{\mathrm{eff}} \approx 366$.  \textbf{Solution:} sample every
$1{,}440$ bars ($n \approx 365$ non-overlapping daily observations per OOS
year) and apply the block bootstrap.

\subsection{The Block Bootstrap}
\label{sec:bootstrap}

The block bootstrap (Künsch, 1989; Liu \& Singh, 1992) provides valid inference
for statistics of dependent time series without parametric assumptions about
the dependence structure.

\begin{definition}[Non-Overlapping Block Bootstrap]
  Given a time series $(X_1, \ldots, X_n)$ and a statistic $T_n$:
  \begin{enumerate}
    \item Choose block length $\ell$ and form $m = \lfloor n/\ell \rfloor$
          non-overlapping blocks $B_j = (X_{(j-1)\ell+1}, \ldots, X_{j\ell})$
          for $j = 1, \ldots, m$.
    \item Draw $m$ blocks $B_{j_1}^*, \ldots, B_{j_m}^*$ uniformly at random
          with replacement from $\{B_1, \ldots, B_m\}$.
    \item Concatenate to form the bootstrap resample
          $X^* = (B_{j_1}^*, \ldots, B_{j_m}^*)$; trim to length $n$.
    \item Compute the bootstrap statistic $T_n^* = T_n(X^*)$.
    \item Repeat $B$ times.  The empirical distribution of
          $\{T_n^{*(b)}\}_{b=1}^B$ approximates the sampling distribution
          of $T_n$.
  \end{enumerate}
\end{definition}

For inference on $\IC$, the block bootstrap preserves temporal autocorrelation
within blocks (capturing weekly seasonality, funding-cycle effects, and
momentum effects) while permuting the dependence structure across blocks.

\paragraph{Block length selection.}
The block length $\ell$ should exceed the effective autocorrelation decay
length of the signal-return interaction.  For the daily non-overlapping series
used here, a block of $\ell = 21$ days (three calendar weeks) captures
relevant dependence: weekly trading patterns, bi-weekly options gamma pinning
cycles near expiry, and multi-week momentum effects.

\paragraph{$p$-value computation.}
For a one-sided test of $H_0 : \IC = 0$ against $H_1 : \IC > 0$:
\begin{equation}
  p = \frac{1}{B}\sum_{b=1}^B \mathbf{1}\!\left[\IC^{*(b)} \geq \widehat{\IC}\right]
  \label{eq:pvalue}
\end{equation}
where $B = 2{,}000$ resamples are used throughout.

\subsection{Pass Criteria (Pre-Committed)}

All criteria were committed before running the IC screen:
\begin{itemize}
  \item \textbf{Unfiltered}: IC/IC* $> 0.5$ and correct sign.
  \item \textbf{DVOL-filtered} ($\geq 54$): ratio $> 1.0$ and bootstrap
        $p \leq 0.05$.
  \item \textbf{Sub-period stability}: IC same-sign in $\geq 3$ of 5
        OOS half-years.
\end{itemize}

IC screen advance triggers a full deep dive with additional gates:
\begin{itemize}
  \item 95\% block-bootstrap CI lower bound strictly positive.
  \item Incremental IC after conditioning on N3$z$ is significant ($p \leq 0.05$).
\end{itemize}

\subsection{Data}

\begin{center}
\begin{tabular}{lll}
\toprule
Dataset & Source & Coverage \\
\midrule
1-minute BTCUSDT klines & Binance FAPI & 2023-01-01 to 2026-05-13 \\
BTC DVOL 1h & Deribit API & 2022-01-01 to 2026-05-15 \\
ETH DVOL 1h & Deribit API & 2022-01-01 to 2026-05-15 (new) \\
Premium index 1m & Binance FAPI & 2023-01-01 to 2026-05-13 \\
Daily liquidations & Binance Data Vision & Unavailable (HTTP 404) \\
\bottomrule
\end{tabular}
\end{center}

The ETH DVOL series was downloaded via the same Deribit endpoint as the
BTC series (\texttt{get\_volatility\_index\_data}, \texttt{currency=ETH}),
yielding 38{,}281 hourly bars.

%% ============================================================
\section{Q4: Mark-Index Basis}
\label{sec:q4}
%% ============================================================

\subsection{Signal Construction}

Both variants derive from the same basis series
$b_t$ (Equation~\ref{eq:basis}), standardised as a 30-day rolling
$z$-score to remove the non-stationary level trend:

\begin{definition}[Basis Signals]
  Let $b_t = (M_t - I_t) / I_t$ be the 1-minute mark-index basis.
  The 30-day rolling mean and standard deviation use 43{,}200 minute-bars
  ($= 30 \times 1440$):
  \begin{align}
    z_t(b) &= \frac{b_t - \mu_{30\mathrm{d}}(b)}{\sigma_{30\mathrm{d}}(b)} \\
    \text{Q4a}_t &= \phantom{-}z_t(b)
      \quad\text{(momentum: positive basis $\to$ long)} \\
    \text{Q4b}_t &= -z_t(b)
      \quad\text{(contrarian: negative basis $\to$ long)}
  \end{align}
  For the daily IC frame, the basis signal at each daily observation is
  the value of $z_t(b)$ at the sampling bar (every 1{,}440 bars).
\end{definition}

\subsection{IC Screen Results}

\begin{center}
\begin{tabular}{lrrrrrrl}
\toprule
Signal & IC & Ratio & $p$ & IC (filt.) & Ratio (filt.) & Stab. & Verdict \\
\midrule
Q4a momentum   & $+0.113$ & $3.79\times$ & $0.000$ & $+0.087$ & $3.39\times$ & $5/5$ & \textbf{KILL} \\
Q4b contrarian & $-0.113$ & $3.79\times$ & $1.000$ & $-0.087$ & $3.39\times$ & $0/5$ & \textbf{KILL} \\
\bottomrule
\end{tabular}
\end{center}

\subsection{Detailed Analysis}

\paragraph{Q4a (momentum).}
The unfiltered IC is strong: $3.79\times$ the maker breakeven, $p \approx 0$,
$5/5$ sub-period stability.  This confirms the momentum hypothesis
(Section~\ref{sec:basis_theory}): a positive basis predicts positive returns
because it signals excess long demand and conviction.

However, the DVOL-filtered bootstrap $p$-value is $0.061$, failing the
pre-committed $0.05$ threshold by a small margin.  The filter restricts
the sample from $n = 863$ to $n = 333$ OOS observations, which substantially
reduces statistical power.  At $n = 333$, the block bootstrap has less
resolution, and a true IC of $+0.087$ requires $p \approx 0.04$--$0.06$
depending on the autocorrelation structure.

\begin{remark}
  Q4a is a borderline case.  Unfiltered, it clearly passes.  DVOL-filtered,
  it misses by $0.011$.  The pre-committed criterion is $p \leq 0.05$;
  $p = 0.061$ fails.  Q4a is formally killed but noted as a candidate
  for re-evaluation as additional OOS data accumulates.  A further year
  of OOS data at the observed IC would move $p$ below $0.05$.
\end{remark}

\paragraph{Q4b (contrarian).}
Q4b is cleanly killed: IC $= -0.113$ (wrong sign), $0/5$ sub-period
stability.  The data unambiguously favour the momentum direction: a
positive basis predicts positive returns, not a negative basis.  The
contrarian hypothesis does not hold in the 2024--2026 OOS period.

\paragraph{Economic interpretation.}
The momentum result (Q4a passes unfiltered) is consistent with the
following structural feature of perpetual markets: the basis is
controlled by the funding mechanism to be approximately zero on average,
but \emph{within} each 8-hour window, a positive basis reflects genuine
excess demand for long exposure.  Market participants who are long despite
a positive basis (and thus paying funding) have a stronger directional
conviction than average --- their presence predicts near-term continuation.

%% ============================================================
\section{Q5: ETH-BTC Cross-DVOL}
\label{sec:q5}
%% ============================================================

\subsection{Signal Construction}

Both $z$-scores use 30-day rolling windows (720 hourly bars):

\begin{definition}[BTC and ETH 30-Day DVOL $z$-Scores]
  \begin{align}
    N3z^{\mathrm{BTC}}_t &= \frac{\DVOL^{\mathrm{BTC}}_t -
      \mu_{720}(\DVOL^{\mathrm{BTC}})}
      {\sigma_{720}(\DVOL^{\mathrm{BTC}})} \\
    N3z^{\mathrm{ETH}}_t &= \frac{\DVOL^{\mathrm{ETH}}_t -
      \mu_{720}(\DVOL^{\mathrm{ETH}})}
      {\sigma_{720}(\DVOL^{\mathrm{ETH}})}
  \end{align}
  where $\mu_{720}$ and $\sigma_{720}$ denote the 720-hour (30-day)
  rolling mean and standard deviation.  Both series are resampled to
  the daily (1{,}440-bar) frame before IC computation.
\end{definition}

\begin{definition}[Cross-DVOL Signals]
  \begin{align}
    \text{Q5a}_t &= N3z^{\mathrm{ETH}}_t
      \quad\text{(ETH $z$-score alone)} \\
    \text{Q5b}_t &= N3z^{\mathrm{BTC}}_t \times N3z^{\mathrm{ETH}}_t
      \quad\text{(product: joint confirmation)} \\
    \text{Q5c}_t &= N3z^{\mathrm{BTC}}_t - N3z^{\mathrm{ETH}}_t
      \quad\text{(difference: BTC-excess)}
  \end{align}
\end{definition}

The correlation between the two $z$-score series provides context for
interpreting each signal:
\[
  \Cor(N3z^{\mathrm{BTC}}, N3z^{\mathrm{ETH}}) = +0.696
\]
This high positive correlation means Q5b (the product) will be positive
whenever both are elevated, and Q5c (the difference) captures the residual
after removing the common factor.

\subsection{IC Screen Results}

\begin{center}
\begin{tabular}{lrrrrrrl}
\toprule
Signal & IC & Ratio & $p$ & IC (filt.) & Ratio (filt.) & Stab. & Verdict \\
\midrule
Q5a ETH $N3z$ alone  & $-0.067$ & $2.24\times$ & $0.975$ & $-0.049$ & $1.93\times$ & $1/5$ & KILL \\
Q5b product          & $-0.015$ & $0.52\times$ & $0.677$ & $-0.007$ & $0.28\times$ & $2/5$ & KILL \\
\textbf{Q5c BTC-excess} & $+0.189$ & $6.35\times$ & $<0.001$ & $+0.210$ & $8.17\times$ & $5/5$ & \textbf{ADVANCE} \\
\bottomrule
\end{tabular}
\end{center}

\subsection{Kill Rationale: Q5a and Q5b}

\paragraph{Q5a (ETH $N3z$ alone).}
IC $= -0.067$: the signal is in the \emph{wrong} direction.  A spike in
ETH implied volatility in isolation predicts \emph{negative} BTC returns,
not positive.  The $p$-value of $0.975$ is a strong one-sided rejection of
the positive-direction hypothesis.  Sub-period stability is $1/5$.

The negative IC is economically interpretable: when ETH spikes above its
own history but BTC does not (i.e., ETH-specific stress), BTC tends to
fall sympathetically via the correlated-asset channel without the compensating
BTC-specific exhaustion mechanism that drives mean reversion.  The
contagion from ETH is more persistent than the reversal.

\paragraph{Q5b (product).}
IC $= -0.015$ ($0.52\times$ ratio): effectively zero.  The product of two
correlated $z$-scores mainly captures the \emph{common} factor.  When both
BTC and ETH are simultaneously elevated (broad crypto stress), the
subsequent BTC return has no systematic directional bias: the broad
crypto fear event is neither reliably followed by a recovery nor a
continuation.  Q5b correctly identifies no edge.

\subsection{Q5c Deep Dive}
\label{sec:q5c_deep}

Q5c advances to a full deep dive.  The complete results from script
\texttt{26\_q5c\_deep\_dive.py}:

\begin{center}
\begin{tabular}{ll}
\toprule
Metric & Value \\
\midrule
OOS IC (non-overlapping daily) & $+0.189$ ($6.35\times$ ratio) \\
Block-bootstrap $p$ & $< 0.0005$ ($B = 2{,}000$ resamples) \\
95\% CI lower bound & $+0.045$ \textbf{(strictly positive; gate passes)} \\
DVOL $\geq 52$ filtered IC & $+0.210$ ($8.17\times$ ratio), $p < 0.001$ \\
Sub-period stability & $5/5$ OOS half-years \\
Incremental IC / IC* (after N3$z$) & $+0.0803$ ($2.70\times$ ratio) \\
$\Cor(\text{Q5c}, N3z^{\mathrm{BTC}})$ & $+0.342$ \\
\midrule
OOS Sharpe (Q5c $\geq 0.5$, DVOL $\geq 52$) & $+3.30$ ($n = 105$ trades) \\
Win rate & $57.1\%$ \\
Mean PnL per trade & $+53.7\bp$ \\
Max drawdown & $-997\bp$ \\
Calmar ratio & $13.58$ \\
\midrule
Combined (Q5c $\geq 0.5$ AND $N3z \geq 0.75$, DVOL $\geq 52$) & $+4.99$ ($n = 64$ trades) \\
\bottomrule
\end{tabular}
\end{center}

\subsubsection{Gate-by-Gate Assessment}

\textbf{Gate 1 --- IC significance.}  $p < 0.0005$ is far below the $0.05$
threshold.  The IC of $+0.189$ is $6.35\times$ the maker breakeven.  This
is one of the strongest IC results in the programme.

\textbf{Gate 2 --- 95\% CI lower bound.}  The block-bootstrap 95\%
confidence interval for the IC has a lower bound of $+0.045 > 0$.  This
means even the pessimistic scenario (5th percentile of the sampling
distribution) implies the true IC is positive and above $\IC^* = 0.035$.
The signal is not marginally significant; it is robustly positive.

\textbf{Gate 3 --- DVOL filter strengthens.}  The IC increases from $+0.189$
(unfiltered) to $+0.210$ (DVOL $\geq 52$), and the ratio from $6.35\times$
to $8.17\times$.  This confirms the signal is concentrated in the same
fear regime that underlies N3 and P3: when the market is already fearful,
BTC-specific excess fear above ETH's level is a stronger contrarian predictor.

\textbf{Gate 4 --- Sub-period stability.}  $5/5$ OOS half-years show
positive IC.  No single 6-month window drives the result.

\textbf{Gate 5 --- Incremental IC.}  After residualising Q5c on
N3$z^{\mathrm{BTC}}$ (removing any component shared with N3), the residual
IC is $+0.0803$ ($2.70\times$ the maker breakeven).  The incremental
predictive power comes from the ETH component ($-N3z^{\mathrm{ETH}}$):
controlling for overall BTC fear, knowing that ETH fear is \emph{lower}
than BTC fear adds substantial directional information.

\subsubsection{Independence from N3}

The correlation $\Cor(\text{Q5c}, N3z^{\mathrm{BTC}}) = +0.342$ indicates
moderate but not high collinearity.  Q5c and N3 are positively correlated
because both include $N3z^{\mathrm{BTC}}$; however, Q5c's additional
$-N3z^{\mathrm{ETH}}$ component breaks the perfect co-movement.

The combined rule (Q5c $\geq 0.5$ AND $N3z \geq 0.75$, DVOL $\geq 52$)
achieves Sharpe $+4.99$ on 64 trades --- substantially better than either
signal alone.  This is consistent with Q5c providing genuine additional
information: when both signals fire, the BTC-specific fear component
has been confirmed by an elevated aggregate fear environment.

\subsubsection{Interpretation}

Q5c = $N3z^{\mathrm{BTC}} - N3z^{\mathrm{ETH}}$ measures
\emph{BTC-specific} fear relative to the crypto baseline.  When this
difference is large and positive:

\begin{enumerate}
  \item BTC's 30-day DVOL has spiked above its own recent history by more
        than ETH's DVOL has spiked above ETH's history.
  \item The excess BTC fear represents forced BTC-specific selling that
        has exhausted: the sellers were BTC-specific (BTC miners, BTC
        ETF arb desks, BTC-collateralised loan liquidations), not
        broad crypto participants.
  \item Broad crypto participants (who hold both BTC and ETH) are less
        dislocated, so their re-entry into BTC is more predictable.
  \item The BTC price recovers toward its fair value relative to ETH
        within 24 hours.
\end{enumerate}

\subsection{Deployment Status}

Q5c was initially deployed as a shadow strategy (\texttt{Q5C\_BTC\_EXCESS},
entry: Q5c $\geq 0.5$ AND DVOL $\geq 52$, hold 24h, maker cost) and
subsequently removed from paper trading when the research programme pivoted
to higher-frequency signal exploration in Phase 5.  The frozen
implementation is preserved in \texttt{strategies/q5c\_btc\_excess.py}
and \texttt{research/active/p4/26\_q5c\_deep\_dive.py}.

%% ============================================================
\section{Q6: Liquidation Exhaustion}
\label{sec:q6}
%% ============================================================

\subsection{Hypothesis}

Forced long-liquidation cascades are a crypto-specific fear mechanism
described in Section~\ref{sec:cascade_theory}.  A spike in daily forced
long-selling marks the final stage of a leveraged sell-off.  Combined
with the DVOL $\geq 54$ filter (ensuring the fear regime is active),
this may identify capitulation bottoms with high precision.

The proposed signal:
\begin{equation}
  \text{Q6}_t = \frac{L_t - \mu_{30}(L)}{\sigma_{30}(L)} \cdot
                \mathbf{1}[\DVOL_t \geq 54]
\end{equation}
where $L_t$ is the total long-liquidation notional (USD) on day $t$.
$\text{Q6}_t > 2$ signals an extreme liquidation event.

\subsection{Data Availability}

The Binance Data Vision liquidation snapshot endpoint:
\begin{center}
\texttt{/data/futures/um/daily/liquidationSnapshot/BTCUSDT/}
\end{center}
returned HTTP 404 for all dates 2022--2026.  The endpoint appears to
have been deprecated or moved; no archived liquidation data was
accessible.

Q6 is formally recorded as \textbf{untestable} in Phase 4.  Alternative
data sources that may enable testing in a future phase:

\begin{itemize}[noitemsep]
  \item \textbf{Coinalyze:} provides historical BTC futures liquidation
        data by exchange, daily granularity, back to 2020.
  \item \textbf{CoinGlass:} provides liquidation heatmaps and historical
        aggregated liquidation data at hourly granularity.
  \item \textbf{Velo Data:} institutional-grade crypto market data
        including liquidation feeds.
\end{itemize}

%% ============================================================
\section{Discussion}
\label{sec:discussion}
%% ============================================================

\subsection{Why Cross-Asset Signals Work When Unilateral Ones Do Not}

The Q5 results demonstrate an important general principle: \emph{conditioning
on a cross-asset spread can produce a stronger signal than the individual
asset levels.}

\begin{itemize}
  \item Q5a (ETH alone, $\IC = -0.067$): fails because ETH-specific fear
        is not a reliable predictor of BTC mean reversion.
  \item Q5b (product, $\IC = -0.015$): fails because broad crypto fear
        (common factor) is not a reliable directional predictor.
  \item Q5c (difference, $\IC = +0.189$): succeeds because it extracts the
        BTC-specific component of fear after removing the systemic ETH-wide
        component.
\end{itemize}

This parallels the logic that made N3 ($z$-score) superior to N1 (DVOL
level, A1 Section 2.4): normalising by a reference removes the common
factor and exposes the idiosyncratic component.

\subsection{The Momentum-Contrarian Tension in Basis Signals}

Q4a (momentum) passes the unfiltered screen but fails the DVOL-filtered
test by $0.011$ in $p$-value.  This represents the coexistence of two
forces:

\begin{itemize}
  \item \textbf{Short-horizon momentum:} within an 8-hour settlement window,
        a positive basis reflects genuine demand and predicts continuation.
  \item \textbf{Medium-horizon reversion:} the funding mechanism ensures
        the basis mean-reverts to zero over multi-day horizons.  At the
        24-hour horizon, both forces are active.
\end{itemize}

The unfiltered IC captures the momentum force in calm periods.  In
high-DVOL periods (the filter), the reversion force is stronger (fear
events cause rapid basis normalisation), and the net directional IC is
weaker.  This regime interaction is why Q4a passes unfiltered but not
filtered.

\subsection{Limitations}

\paragraph{DVOL-filtered sample size.}
The DVOL $\geq 54$ filter reduces the OOS sample from 863 to 333
observations.  At $n = 333$ with block size 21, the block bootstrap
has only $333/21 \approx 16$ independent blocks, limiting the
precision of the $p$-value estimate.  The borderline Q4a result
($p = 0.061$) should be interpreted as ``near the threshold'' rather
than as a qualitative failure.

\paragraph{Q6 data gap.}
The inability to test Q6 is the most significant limitation of Phase 4.
The liquidation exhaustion hypothesis is economically well-motivated and
empirically untested in this programme.  If daily liquidation data becomes
available, Q6 should be the first signal in the Phase 6 queue.

\paragraph{Correlation regime stability.}
The BTC-ETH DVOL correlation of $+0.696$ is measured over the 2022--2026
sample.  During structural regime changes (e.g., a period of sustained
ETH-specific risk such as a major DeFi collapse, or a BTC-specific shock
such as a mining regulatory crisis), the correlation may shift, affecting
the Q5c signal's effective information content.  No regime stability test
was performed in Phase 4.

%% ============================================================
\section{Summary and Kill Log}
\label{sec:summary}
%% ============================================================

\begin{center}
\begin{tabular}{lll}
\toprule
Signal & Verdict & Primary reason \\
\midrule
Q4a Basis momentum   & Kill & DVOL-filtered $p = 0.061 > 0.05$ (borderline; candidate for future re-test) \\
Q4b Basis contrarian & Kill & Wrong direction ($\IC = -0.113$, $0/5$ stability) \\
Q5a ETH $N3z$ alone  & Kill & Wrong direction in OOS ($\IC = -0.067$, $1/5$ stability) \\
Q5b BTC$\times$ETH   & Kill & Near-zero IC ($0.52\times$), $2/5$ stability \\
Q5c BTC-excess $N3z$ & \textbf{Advance} (superseded) & All gates pass; $5/5$; Sharpe $+3.30$; CI lower $+0.045$ \\
Q6 Liquidation Exhaustion & Untestable & Binance Data Vision endpoint unavailable \\
\bottomrule
\end{tabular}
\end{center}

Research scripts: \texttt{data/download\_phase4.py},
\texttt{research/active/p4/25\_phase4\_ic\_screen.py},
\texttt{research/active/p4/26\_q5c\_deep\_dive.py}.

%% ============================================================
\section{Conclusion}
\label{sec:conclusion}
%% ============================================================

Phase 4 tested three new signal families and validated one: Q5c, the
BTC-excess implied volatility signal.

The key results:

\begin{enumerate}
  \item \textbf{Q5c is a robust, independent signal.}
        IC $= +0.189$ ($6.35\times$ maker breakeven), block-bootstrap
        $p < 0.0005$, 95\% CI lower bound $+0.045$ (strictly positive),
        $5/5$ sub-period stability, and incremental IC $2.70\times$
        after N3$z$ control.  OOS Sharpe $+3.30$ on 105 trades.

  \item \textbf{The cross-asset spread (Q5c) outperforms its components.}
        ETH $N3z$ alone (Q5a) has IC $= -0.067$.  BTC$\times$ETH product
        (Q5b) has IC $= -0.015$.  The difference Q5c has IC $= +0.189$.
        Conditioning on BTC-specific excess fear, after removing the
        systemic ETH component, is the key filtering step.

  \item \textbf{The basis is a momentum signal but fails the DVOL-filtered
        gate.}  Q4a confirms the momentum hypothesis (positive basis
        predicts positive returns) in calm periods, but the DVOL-filtered
        $p = 0.061$ fails the pre-committed $0.05$ threshold.

  \item \textbf{Q6 remains a high-priority untested hypothesis.}
        The liquidation exhaustion mechanism is economically compelling
        (Section~\ref{sec:cascade_theory}) but data unavailability
        prevents testing.  Sourcing alternative liquidation data is a
        Phase 6 priority.
\end{enumerate}

Q5c was deployed to shadow paper trading and subsequently superseded when
the research programme pivoted to higher-frequency signals in Phase 5
(A4).  Its research trail is preserved in \texttt{research/active/p4/}.

\vspace{1em}
\noindent\textit{Data: Binance FAPI 1m klines (2023--2026); Deribit
BTC and ETH DVOL 1h (2022--2026); Binance FAPI premium index 1m
(2023--2026).}

%% ============================================================
\section*{References}
%% ============================================================
\begingroup
\setlength{\parindent}{-2em}
\setlength{\leftskip}{2em}
\setlength{\parskip}{6pt}
\small

Britten-Jones, M., \& Neuberger, A. (2000). Option prices, implied price processes, and stochastic volatility. \textit{The Journal of Finance}, \textit{55}(2), 839--866. \url{https://doi.org/10.1111/0022-1082.00228}

Diebold, F. X., \& Yilmaz, K. (2012). Better to give than to receive: Predictive directional measurement of volatility spillovers. \textit{International Journal of Forecasting}, \textit{28}(1), 57--66. \url{https://doi.org/10.1016/j.ijforecast.2011.02.006}

Künsch, H. R. (1989). The jackknife and the bootstrap for general stationary observations. \textit{The Annals of Statistics}, \textit{17}(3), 1217--1241. \url{https://doi.org/10.1214/aos/1176347265}

Liu, R. Y., \& Singh, K. (1992). Moving blocks jackknife and bootstrap capture weak convergence. In R. LePage \& L. Billard (Eds.), \textit{Exploring the Limits of Bootstrap} (pp. 225--248). Wiley.

\endgroup

\end{document}
